% ============================================================================
% PREÂMBULO — "Full Stack Moderno: do zero ao sênior"
% Compilação: XeLaTeX + biber (ver scripts/build-pdf.sh e scripts/build-epub.sh)
% Requisitos: TeX Live 2023+ (ou TinyTeX) com os pacotes listados no README.
% ============================================================================

% ---------- Página e geometria (digital: A4, margens confortáveis) ----------
\usepackage[a4paper, top=2.4cm, bottom=2.6cm, inner=2.6cm, outer=2.2cm,
            headheight=16pt, footskip=1.2cm, marginparwidth=0cm]{geometry}

% ---------- Fontes (XeLaTeX, Unicode seguro) ----------
\usepackage{fontspec}
% Fontes carregadas por NOME DE ARQUIVO via kpathsea (independente do
% fontconfig do sistema operacional): TeX Gyre com fallback para Latin Modern.
\IfFileExists{texgyrepagella-regular.otf}{%
  \setmainfont{TeX Gyre Pagella}[
    UprightFont = texgyrepagella-regular.otf,
    BoldFont = texgyrepagella-bold.otf,
    ItalicFont = texgyrepagella-italic.otf,
    BoldItalicFont = texgyrepagella-bolditalic.otf,
  ]%
}{%
  \setmainfont{Latin Modern Roman}[
    UprightFont = lmroman10-regular.otf,
    BoldFont = lmroman10-bold.otf,
    ItalicFont = lmroman10-italic.otf,
    BoldItalicFont = lmroman10-bolditalic.otf,
  ]%
}
\IfFileExists{texgyreheros-regular.otf}{%
  \setsansfont{TeX Gyre Heros}[
    UprightFont = texgyreheros-regular.otf,
    BoldFont = texgyreheros-bold.otf,
    ItalicFont = texgyreheros-italic.otf,
    BoldItalicFont = texgyreheros-bolditalic.otf,
  ]%
}{%
  \setsansfont{Latin Modern Sans}[
    UprightFont = lmsans10-regular.otf,
    BoldFont = lmsans10-bold.otf,
    ItalicFont = lmsans10-oblique.otf,
    BoldItalicFont = lmsans10-boldoblique.otf,
  ]%
}
% Fonte monoespaçada: Inconsolata se disponível; senão Latin Modern Mono.
\IfFileExists{Inconsolata-Regular.otf}{%
  \setmonofont{Inconsolata}[
    Scale = 0.9,
    UprightFont = Inconsolata-Regular.otf,
    BoldFont = Inconsolata-Bold.otf,
  ]%
}{%
  \setmonofont{Latin Modern Mono}[
    Scale = 0.9,
    UprightFont = lmmono10-regular.otf,
    BoldFont = lmmono10-regular.otf,   % LM Mono não tem bold dedicado
    ItalicFont = lmmono10-italic.otf,
    BoldItalicFont = lmmono10-italic.otf,
  ]%
}

% ---------- Idioma e hifenização (pt-BR) ----------
\usepackage[brazilian]{babel}
\shorthandoff{"}   % o babel português ativa `"`; desligamos para permitir aspas retas em \texttt{...}

% ---------- Tipografia fina ----------
\usepackage{microtype}
\usepackage{amsmath, amssymb}
\emergencystretch=3em          % evita overfull hbox em parágrafos densos
\setlength{\parskip}{0.15em}

% ---------- Cores (paleta editorial) ----------
\usepackage[dvipsnames, table]{xcolor}
\definecolor{primaria}{HTML}{1F4E79}   % azul-escuro institucional
\definecolor{secundaria}{HTML}{2E86AB} % azul-médio
\definecolor{acento}{HTML}{C0392B}     % vermelho de alerta
\definecolor{verde}{HTML}{1E8449}
\definecolor{verdeclaro}{HTML}{E8F8F0}
\definecolor{laranja}{HTML}{B9770E}
\definecolor{laranjacla}{HTML}{FCF3E7}
\definecolor{azulclaro}{HTML}{EAF2F8}
\definecolor{cinzaclaro}{HTML}{F2F3F4}
\definecolor{codigo-fundo}{HTML}{FAFAFA}
\definecolor{codigo-borda}{HTML}{CCCCCC}
\definecolor{comentario}{HTML}{5F7A61}
\definecolor{stringcor}{HTML}{9A4D1E}
\definecolor{keywordcor}{HTML}{1F4E79}

% ---------- Listagens de código (à prova de estouro de borda) ----------
\usepackage{listings}

% Linguagens definidas localmente (garantia de funcionamento em qualquer TeX Live)
\lstdefinelanguage{JavaScript}{
  morekeywords={const,let,var,function,return,if,else,for,while,do,switch,case,
    break,continue,new,class,extends,super,this,typeof,instanceof,delete,void,
    await,async,try,catch,finally,throw,import,export,from,default,yield,of,in,
    static,get,set,true,false,null,undefined,NaN,Infinity},
  sensitive=true,
  morecomment=[l]{//},
  morecomment=[s]{/*}{*/},
  morestring=[b]',
  morestring=[b]"
}
\lstdefinelanguage{TypeScript}{
  language=JavaScript,
  morekeywords={interface,type,enum,namespace,declare,readonly,implements,
    public,private,protected,abstract,keyof,infer,as,is,unknown,never,any,
    string,number,boolean,symbol,bigint,object,Record,Partial,Required,Readonly,
    Pick,Omit,Exclude,Extract,ReturnType,Parameters,Promise,Array}
}
\lstdefinelanguage{Bash}{
  morekeywords={echo,cd,ls,mkdir,rm,mv,cp,cat,curl,npm,npx,node,git,docker,
    docker-compose,export,source,if,then,else,fi,for,in,do,done,function,local,
    case,esac,while,test,pwd,touch,grep,sudo,chmod,chown,ps,kill},
  sensitive=false,
  morecomment=[l]{\#},
  morestring=[b]'
}
\lstdefinelanguage{SQL}{
  morekeywords={SELECT,FROM,WHERE,INSERT,INTO,VALUES,UPDATE,SET,DELETE,CREATE,
    TABLE,ALTER,DROP,INDEX,PRIMARY,KEY,FOREIGN,REFERENCES,JOIN,INNER,LEFT,RIGHT,
    FULL,ON,GROUP,BY,ORDER,HAVING,LIMIT,OFFSET,AS,DISTINCT,AND,OR,NOT,NULL,IS,
    BETWEEN,LIKE,IN,EXISTS,UNION,ALL,CASE,WHEN,THEN,ELSE,END,BEGIN,COMMIT,
    ROLLBACK,TRANSACTION,VARCHAR,TEXT,INTEGER,BIGINT,SERIAL,UUID,BOOLEAN,TIMESTAMP,
    DATE,NUMERIC,DECIMAL,CHECK,DEFAULT,UNIQUE},
  sensitive=true,
  morecomment=[l]{--}
}
\lstdefinelanguage{HTML}{
  morekeywords={html,head,body,title,meta,link,script,style,div,span,p,a,h1,h2,h3,
    h4,h5,h6,ul,ol,li,table,thead,tbody,tr,th,td,form,input,label,button,select,
    option,textarea,img,section,article,header,footer,nav,main,aside,figure,
    figcaption,video,audio,canvas,iframe,strong,em,code,pre,blockquote,br,hr,
    small,mark,time,address,details,summary},
  sensitive=true,
  morecomment=[s]{<!--}{-->}
}
\lstdefinelanguage{CSS}{
  morekeywords={color,background,background-color,font,font-size,font-family,
    font-weight,margin,padding,border,border-radius,display,position,top,right,
    bottom,left,width,height,min-width,max-width,flex,flex-direction,grid,
    gap,align-items,align-self,justify-content,text-align,text-decoration,
    line-height,opacity,z-index,overflow,box-sizing,box-shadow,transition,
    transform,animation,media,import,root,keyframes,content,float,clear,cursor,
    list-style,outline,visibility,vertical-align,white-space,word-break},
  sensitive=true,
  morecomment=[s]{/*}{*/}
}
\lstdefinelanguage{JSON}{
  morekeywords={true,false,null},
  sensitive=true
}

\lstdefinestyle{livro}{
  basicstyle=\ttfamily\footnotesize,
  keywordstyle=\color{keywordcor}\bfseries,
  commentstyle=\color{comentario}\itshape,
  stringstyle=\color{stringcor},
  showstringspaces=false,
  breaklines=true,            % quebra linhas longas (evita estouro)
  breakatwhitespace=true,
  postbreak=\mbox{\textcolor{secundaria}{\ensuremath{\hookrightarrow}}\space},
  columns=fullflexible,       % espaçamento flexível (evita estouro)
  keepspaces=true,
  frame=single,
  rulecolor=\color{codigo-borda},
  backgroundcolor=\color{codigo-fundo},
  numbers=left,
  numberstyle=\tiny\color{gray!70},
  stepnumber=1,
  xleftmargin=1.6em,
  xrightmargin=0.4em,
  aboveskip=0.9em,
  belowskip=0.9em,
  captionpos=b,
  tabsize=2,
  showtabs=false,
  literate= % acentuação pt-BR dentro do código (comentários e strings)
    {á}{{\'a}}1 {à}{{\`a}}1 {â}{{\^a}}1 {ã}{{\~a}}1 {ä}{{\"a}}1
    {é}{{\'e}}1 {è}{{\`e}}1 {ê}{{\^e}}1 {ë}{{\"e}}1
    {í}{{\'i}}1 {ì}{{\`i}}1 {î}{{\^i}}1 {ï}{{\"i}}1
    {ó}{{\'o}}1 {ò}{{\`o}}1 {ô}{{\^o}}1 {õ}{{\~o}}1 {ö}{{\"o}}1
    {ú}{{\'u}}1 {ù}{{\`u}}1 {û}{{\^u}}1 {ü}{{\"u}}1
    {ç}{{\c{c}}}1 {Ç}{{\c{C}}}1
    {Á}{{\'A}}1 {À}{{\`A}}1 {Â}{{\^A}}1 {Ã}{{\~A}}1 {Ä}{{\"A}}1
    {É}{{\'E}}1 {È}{{\`E}}1 {Ê}{{\^E}}1 {Ë}{{\"E}}1
    {Í}{{\'I}}1 {Ì}{{\`I}}1 {Î}{{\^I}}1 {Ï}{{\"I}}1
    {Ó}{{\'O}}1 {Ò}{{\`O}}1 {Ô}{{\^O}}1 {Õ}{{\~O}}1 {Ö}{{\"O}}1
    {Ú}{{\'U}}1 {Ù}{{\`U}}1 {Û}{{\^U}}1 {Ü}{{\"U}}1
    {º}{{\textordmasculine}}1 {ª}{{\textordfeminine}}1
    {—}{{---}}1 {–}{{--}}1 {…}{{\ldots}}1
}

% ---------- Caixas didáticas (quebráveis, sem estouro) ----------
\usepackage[most]{tcolorbox}
\newtcolorbox{caixadica}{breakable, enhanced, arc=1.5mm, boxrule=0.6pt,
  colback=verdeclaro, colframe=verde, coltitle=white,
  title={\small\bfseries Dica}, fonttitle=\bfseries,
  left=2.5mm, right=2.5mm, top=1.6mm, bottom=1.6mm}
\newtcolorbox{caixaarmadilha}{breakable, enhanced, arc=1.5mm, boxrule=0.6pt,
  colback=laranjacla, colframe=laranja, coltitle=white,
  title={\small\bfseries Armadilha comum}, fonttitle=\bfseries,
  left=2.5mm, right=2.5mm, top=1.6mm, bottom=1.6mm}
\newtcolorbox{caixaprojeto}{breakable, enhanced, arc=1.5mm, boxrule=0.8pt,
  colback=azulclaro, colframe=primaria, coltitle=white,
  title={\small\bfseries Projeto do capítulo}, fonttitle=\bfseries,
  left=2.5mm, right=2.5mm, top=1.6mm, bottom=1.6mm}
\newtcolorbox{caixaconceito}{breakable, enhanced, arc=1.5mm, boxrule=0.5pt,
  colback=cinzaclaro, colframe=gray!60, coltitle=black,
  title={\small\bfseries Conceito-chave}, fonttitle=\bfseries,
  left=2.5mm, right=2.5mm, top=1.6mm, bottom=1.6mm}

% ---------- Cabeçalhos e rodapés ----------
\usepackage{fancyhdr}
\pagestyle{fancy}
\fancyhf{}
\fancyhead[L]{\small\color{gray!70}\leftmark}
\fancyhead[R]{\small\color{primaria}\thepage}
\renewcommand{\headrulewidth}{0.4pt}
\renewcommand{\headrule}{\color{secundaria}\hrule width\headwidth height 0.4pt}
\fancypagestyle{plain}{%
  \fancyhf{}
  \fancyhead[R]{\small\color{primaria}\thepage}
  \renewcommand{\headrulewidth}{0pt}}

% ---------- Formatação de capítulos e seções ----------
\usepackage{titlesec}
\titleformat{\chapter}[display]
  {\normalfont\sffamily\bfseries\color{primaria}}
  {\filleft\Large\MakeUppercase{\chaptertitlename} \thechapter}
  {0.4em}
  {\Huge}
  [\vspace{0.4em}{\color{secundaria}\titlerule[1pt]}]
\titlespacing*{\chapter}{0pt}{-30pt}{26pt}
\titleformat{\section}{\normalfont\Large\bfseries\color{primaria}}{\thesection}{0.7em}{}
\titleformat{\subsection}{\normalfont\large\bfseries\color{secundaria}}{\thesubsection}{0.7em}{}
\titleformat{\subsubsection}{\normalfont\normalsize\bfseries\color{black}}{\thesubsubsection}{0.7em}{}

% ---------- Sumário: folga para números de seção longos (ex.: 12.15) ----------
% O box padrão de 2.3em deixa apenas ~0.7pt entre o número de 4 dígitos e o
% título, parecendo colado. 2.6em garante ~4.4pt de respiro para o pior caso.
\makeatletter
\renewcommand*\l@section{\@dottedtocline{1}{1.5em}{2.6em}}
\makeatother

% ---------- Tabelas e legendas ----------
\usepackage{booktabs}
\usepackage{longtable}
\usepackage{tabularx}
\usepackage{caption}
\captionsetup{font=small, labelfont={bf,color=primaria}, skip=6pt}
\captionsetup[lstlisting]{font=small, labelfont={bf,color=primaria}}

% ---------- Listas e ambientes didáticos ----------
\usepackage{enumitem}
\newlist{objetivos}{itemize}{1}
\setlist[objetivos]{label=\textcolor{secundaria}{\textbullet},
  leftmargin=*, itemsep=0.3em, topsep=0.5em}
\newlist{exercicios}{enumerate}{1}
\setlist[exercicios]{label=\textbf{Exercício \arabic*.},
  leftmargin=*, itemsep=0.7em, topsep=0.5em, before={\smallskip}}
\newlist{desafios}{enumerate}{1}
\setlist[desafios]{label=\textbf{Desafio \arabic*.},
  leftmargin=*, itemsep=0.7em, topsep=0.5em, before={\smallskip}}
\newlist{seniores}{enumerate}{1}
\setlist[seniores]{label=\textbf{Sênior \arabic*.},
  leftmargin=*, itemsep=0.7em, topsep=0.5em, before={\smallskip}}
\newlist{passos}{enumerate}{1}
\setlist[passos]{label=\textbf{Passo \arabic*.},
  leftmargin=*, itemsep=0.6em, topsep=0.4em, before={\smallskip}}

% ---------- Índice remissivo (makeindex, estilo em indice.ist) ----------
\usepackage[original]{imakeidx}
\makeindex[options={-s indice.ist}, title={Índice Remissivo}, intoc]

% ---------- Hiperlinks e URLs (xurl evita estouro em URLs longas) ----------
\usepackage[hidelinks, bookmarksnumbered=true,
  pdftitle={Full Stack Moderno: do zero ao sênior},
  pdfauthor={Autor do Livro},
  pdfsubject={Desenvolvimento Web Full Stack},
  pdfkeywords={full stack, TypeScript, React, Next.js, PostgreSQL, Docker, AWS},
  pdfcreator={XeLaTeX}]{hyperref}
\usepackage{xurl}

% ---------- Referências bibliográficas (BibLaTeX + biber) ----------
\usepackage[style=numeric, backend=biber, sortcites=true]{biblatex}
\addbibresource{referencias.bib}

% ---------- Campos editáveis do livro (veja README.md) ----------
\newcommand{\autordolivro}{[Nome do Autor]}
\newcommand{\versaodolivro}{0.2}
\newcommand{\datadolivro}{Agosto de 2026}

% ---------- Diversos ----------
\usepackage{etoolbox}
\usepackage{tikz}
\usetikzlibrary{arrows.meta, positioning, shapes.geometric}
% Diagramas vetoriais: TikZ no PDF; nota textual no EPUB (o pipeline tex4ebook
% não converte TikZ sem dvisvgm). O diagrama em si nunca é processado no EPUB.
\newcommand{\diagramatikz}[1]{%
  \ifdefined\HCode
    \begin{minipage}{0.92\linewidth}\centering
    \footnotesize\textit{Diagrama vetorial — disponível na versão PDF deste
    livro.}
    \end{minipage}%
  \else
    #1%
  \fi
}
\newcommand{\codigopath}{../codigo}   % caminho base do repositório de código
